\chapter{Metodologia}
\label{cap:metodo}

A pesquisa compreende desenvolvimento, verificação funcional, calibração dos instrumentos e comparação experimental. A verificação do software identifica se cada mecanismo funciona nas condições testadas; a comparação exige o conjunto de repetições e controles especificado neste capítulo. Os registros da primeira etapa não foram utilizados como substitutos da coleta definitiva.

\section{Tipo de pesquisa}

\subsection{Natureza da pesquisa}
A pesquisa é aplicada porque constrói um protótipo para apoiar decisões de engenharia de software. O procedimento principal é experimental: duas configurações de comunicação foram submetidas a cargas e falhas controladas. Desempenho e recuperação dependem de condições de execução e, portanto, não podem ser estabelecidos apenas pela descrição das tecnologias. A definição prévia de variáveis e critérios de análise segue as recomendações de \citeonline{kitchenham2002}.

\subsection{Abordagem metodológica}
A abordagem é quantitativa. Latência, throughput, erros, recursos, backlog e recuperação são expressos por medidas observáveis. A interpretação relaciona os números aos mecanismos executados e às limitações do ambiente. A unidade experimental é a execução planejada, com elegibilidade por métrica explicitada na revisão analítica; os eventos de uma mesma execução não são repetições experimentais independentes. O Quadro~\ref{qua:termos} define os termos adotados; os identificadores técnicos ficam no Apêndice~\ref{ap:rastreabilidade}.

\begin{quadro}[htbp]
\centering\small
\caption{Termos da organização experimental}\label{qua:termos}
\begin{tabular}{|p{0.18\textwidth}|p{0.71\textwidth}|}
\hline
Termo & Significado neste trabalho\\ \hline
Evento & Um registro sintético enviado à API. Possui identificador único; milhares de eventos podem pertencer à mesma coleta.\\ \hline
Coleta & Uma execução de uma variante em uma condição, com aquecimento, medição e observação posterior. É a unidade experimental, não um evento.\\ \hline
Condição & Combinação de perfil de carga ou falha, taxa e variante. Foram treze condições com cinco coletas cada.\\ \hline
Par & Duas coletas, uma por variante, com o mesmo conjunto determinístico de eventos na mesma condição comparável.\\ \hline
Bloco & Agrupamento de uma repetição de todas as condições, usado para organizar e contrabalançar a ordem. Foram cinco blocos.\\ \hline
Piloto & Verificação prévia dos instrumentos e parâmetros. Não integra as repetições definitivas.\\ \hline
Revisão analítica & Emenda posterior ao início da campanha que definiu elegibilidade por indicador diante de uma evidência temporal incompleta, sem alterar os dados ou repetir a coleta.\\ \hline
\end{tabular}
\fonte{Elaborado pelo autor (2026).}
\end{quadro}


A relação com os objetivos específicos organiza o procedimento. Contrato e regras comuns estabelecem equivalência funcional; APIs, consumidor e repositório materializam as variantes; carga, instrumentação e injeção de falha operacionalizam os cenários; consolidação e análise dos dados permitem discutir os trade-offs. A equivalência funcional controla a regra de processamento, mas não elimina diferenças de topologia: Kafka e consumidor adicionais integram o custo observado da variante assíncrona.

\subsection{Procedimentos técnicos}
Os procedimentos incluem estudo de artigos científicos, desenvolvimento incremental, testes unitários e de integração, calibração e experimento controlado. O levantamento bibliográfico fundamenta conceitos e situa os estudos relacionados; a resposta à pergunta de pesquisa foi construída a partir das medições no protótipo.

\section{Procedimentos metodológicos}

\subsection{Métodos de coleta de informações}
A fundamentação utiliza artigos científicos publicados em periódicos e anais de conferências ou workshops. Documentações oficiais apoiam a descrição das ferramentas, separadamente dos estudos científicos. Os tipos de contribuição e as limitações dos estudos são discutidos no Capítulo~\ref{cap:revisao}.

Na coleta computacional, cada evento possui UUID e instante de ocorrência definidos antes do envio. O conjunto completo de campos é preservado em reentregas. Um identificador UUID de execução, \texttt{run\_id}, é enviado no cabeçalho \texttt{X-Run-ID} e propagado ao consumidor por cabeçalho Kafka. Esse metadado identifica a coleta, sem alterar o conteúdo usado no hash.

O gerador de carga com k6 e o analisador independente em Java foram implementados no diretório de experimentação do projeto. A calibração inicial foi registrada em 26 de outubro de 2026; em 27 de outubro, a revisão reduzida R2 passou nas nove validações adicionais e foi congelada antes de qualquer coleta definitiva. Esses instrumentos não são módulos das APIs nem substituem seu processamento. O executor registra fases separadas de aquecimento e medição e utiliza um projeto Docker Compose isolado. As aplicações Java emitem marcos JSON de etapa, UUID, execução, resultado e tempo. O marco \texttt{commit\_completed} é emitido pelo repositório comum depois do retorno da transação, tanto na variante síncrona quanto no consumidor. Uma repetição idempotente é identificada como \texttt{duplicate} e não representa uma nova linha persistida.

A coleta mantém três instantes distintos: ocorrência do fato, conclusão observada pela aplicação depois do commit e recebimento da resposta HTTP. O campo \texttt{persisted\_at} é preenchido durante a inserção no PostgreSQL; não representa o término exato do commit. A consulta \texttt{GET /audit/\{event\_id\}} verifica a visibilidade da linha, mas seu instante de resposta também não substitui o marco pós-commit.

O tempo de requisição é observado pelo gerador de carga. Na correlação entre processos, utiliza-se uma referência UTC comum; resolução e consistência foram verificadas no piloto. A presença de timestamps em nanossegundos nos logs não aumenta por si só a resolução da medida de ponta a ponta. Relógios monotônicos são usados nas durações locais. Valores incompatíveis são sinalizados como falha de instrumentação, sem truncamento silencioso; essas verificações não comprovam ausência de todo desvio entre relógios.

Na execução experimental, estatísticas Docker registram CPU e memória por contêiner. Offsets do tópico e do grupo permitem observar o lag; consultas ao PostgreSQL verificam existência, hash e unicidade. O executor preserva arquivos brutos, configuração, versões, identificação das imagens e checksums. Esses instrumentos já foram implementados e validados no piloto; a campanha definitiva, distinta das verificações funcionais da Seção~\ref{sec:implementacao-real}, fornece as repetições necessárias à análise comparativa.

\subsection{Métodos de análise das informações}
A análise parte dos dados brutos preservados, correlacionados por execução e UUID. Para cada repetição são calculados quantidade oferecida, aceita e concluída; média, mediana, desvio-padrão e percentis p95 e p99 das latências; throughput concluído; erros por categoria; backlog máximo; reentregas e recuperação. A taxa oferecida efetiva e as iterações que o gerador não conseguiu iniciar são apresentadas separadamente. Na consolidação HTTP válido, percentis são calculados por posto mais próximo sobre as amostras de cada execução; preservam-se também os resumos originais do analisador e do k6, cujos métodos de interpolação podem produzir pequenas diferenças. Não se misturam suas populações ou percentis.

Primeiro são examinados os resultados de cada repetição. Para o orçamento reduzido de cinco repetições por condição, a análise predeclarada é descritiva: valores individuais, média, mediana, desvio-padrão amostral, mínimo, máximo e diferenças pareadas entre variantes. A direção das cinco diferenças é discutida sem convertê-la em teste de superioridade ou equivalência estatística. O tamanho foi definido pelo orçamento operacional, sem cálculo de poder, e a análise descritiva foi fixada antes da campanha. Com apenas cinco pares, até cinco diferenças concordantes teriam probabilidade bilateral de 6,25\% em um teste exato de sinais sob a hipótese de direções equiprováveis e pares independentes; o cálculo ilustra a baixa resolução inferencial, não constitui teste acrescentado aos resultados. A dependência possível do ambiente compartilhado reforça a cautela. Não foram acrescentadas execuções para obter significância ou favorecer uma hipótese. Testes com poucos eventos destinados à depuração não sustentam inferências sobre desempenho ou superioridade arquitetural.

\section{Metodologia de desenvolvimento do sistema}

\subsection{Modelo de desenvolvimento de software}
O desenvolvimento é incremental: contrato, serialização e persistência; variante síncrona; produtor e consumidor; tratamento de falhas; instrumentação; calibração; coleta definitiva. Cada incremento mantém requisitos, código, diagramas e verificações correspondentes. Casos com doubles de bibliotecas verificam decisões locais; testes de integração utilizam Kafka e PostgreSQL reais.

\subsection{Tecnologias utilizadas}
O protótipo utiliza Java com bytecode alvo 21, Spring Boot 3.5.16, Spring JDBC, Flyway, PostgreSQL 16, Apache Kafka 3.9.1 e Docker Compose. O contrato \texttt{AuditEvent} e o repositório comuns reduzem divergências funcionais. PostgreSQL oferece transações e unicidade; Kafka oferece retenção e posições de consumo; Compose descreve a implantação local. O gerador k6, o analisador Java, a conciliação e os coletores foram implementados e calibrados. As imagens e o código instrumentador dessa calibração foram preservados na revisão reduzida; não foram reconstruídos durante as medições.

Cada coleta registrou sistema operacional, processador, memória, versões efetivas, identificadores das imagens, limites de CPU e memória, partições, replicação, parâmetros de timeout e retry e propriedades da carga. O commit Git da implementação Java e checksums dos arquivos compuseram o manifesto. Credenciais permaneceram na configuração privada do ambiente e não integraram relatórios, capturas ou dados experimentais.

\subsection{Ferramentas de geração de carga}
Foram examinadas as documentações oficiais do Apache JMeter e do Grafana k6. O Quadro~\ref{qua:ferramentas} registra a comparação funcional utilizada na escolha do instrumento. Não foi executado um benchmark entre os geradores e, portanto, a escolha não afirma que k6 tenha melhor desempenho que JMeter.

\begin{quadro}[htbp]
\centering
\caption{Comparação dos instrumentos de geração de carga}
\label{qua:ferramentas}
\small
\begin{tabular}{|p{0.19\textwidth}|p{0.35\textwidth}|p{0.35\textwidth}|}
\hline
\textbf{Critério} & \textbf{Apache JMeter} & \textbf{Grafana k6} \\ \hline
Plano & Plano de teste JMX, interface gráfica e propriedades parametrizáveis. & Funções e cenários definidos em JavaScript. \\ \hline
Execução & Modo de linha de comando recomendado para carga; listeners gráficos aumentam o consumo. & Linha de comando e execução em contêiner. \\ \hline
Chegada aberta & Open Model Thread Group, classificado na documentação como experimental. & Executores de taxa constante e variável, independentes da duração das respostas. \\ \hline
Dados de saída & JTL/CSV; recomenda-se limitar os dados coletados. & JSON por linha, comprimível, com tags de correlação. \\ \hline
\end{tabular}
\fonte{Elaborado pelo autor com base nas documentações oficiais das ferramentas, consultadas em 30 out. 2026.}
\end{quadro}

As boas práticas do JMeter recomendam carga em linha de comando e limitar listeners gráficos \cite{jmeterdocs}. O k6 oferece executores de chegada independentes do tempo de resposta e saída JSON por amostra \cite{k6docs,k6json}. A existência de um modelo aberto também no JMeter impede atribuir essa capacidade exclusivamente ao k6. Essas fontes documentam instrumentos; não substituem os artigos científicos da revisão.

O k6 foi escolhido pela integração com o protocolo já implementado: scripts versionados permitem gerar os UUIDs e os conteúdos sintéticos, separar aquecimento e medição, registrar a fase de cada amostra e enviar o identificador da execução no cabeçalho HTTP. A imagem utilizada é a versão 2.3.0, fixada por digest. O analisador Java cruza os dados do gerador com os marcos de transação e os registros exportados do PostgreSQL. A exportação por evento e a instrumentação também consomem recursos, razão pela qual precisam permanecer constantes nas comparações e ser reconhecidas como limitação.


As documentações oficiais de Kafka e PostgreSQL delimitam as garantias de entrega e de isolamento usadas na implementação \cite{kafkadocs,postgresdocs}. O coletor de recursos utiliza as estatísticas de contêineres documentadas pelo Docker \cite{dockerdocs}; não mede todo o computador e não interpreta contêiner parado como consumo zero saudável.

\subsection{Arquitetura do sistema}
As variantes compartilham \texttt{AuditEvent}, serialização determinística, SHA-256 e repositório PostgreSQL. Na síncrona, a resposta 201 sucede a transação ou o reconhecimento de duplicata. Na assíncrona, a API aguarda confirmação individual de publicação com \texttt{acks=all}; o consumidor processa posteriormente. A resposta 202 confirma essa publicação, sem atestar persistência no PostgreSQL. Um timeout de publicação é classificado como resultado incerto, pois a mensagem pode ter sido entregue sem que a confirmação chegasse ao cliente.

O consumidor desabilita commit automático de offsets. Após inserção ou reentrega idêntica, confirma a posição consumida. Erros transitórios de persistência recebem tentativas limitadas com espera exponencial; mensagens inválidas, conflitos ou tentativas esgotadas são encaminhados à DLQ. A posição de origem somente é confirmada depois do ACK da DLQ. Se a publicação na DLQ ou a confirmação da posição falhar, a sessão é retomada sem avançar além da mensagem não resolvida. Isso permite reentrega e exige idempotência; não estabelece uma transação única entre Kafka e PostgreSQL.

A topologia local contém um broker e uma réplica por tópico. A confirmação configurada não demonstra tolerância à perda física do broker nem alta disponibilidade de cluster. Limites de recursos do ambiente de desenvolvimento estão explícitos no Compose; sua adequação às configurações definitivas foi verificada no piloto.

\section{Procedimentos de teste e validação}

\subsection{Tipos de testes aplicados}
Testes unitários Java verificam contrato, serialização, hash, decisões de confirmação, retry, DLQ e respostas HTTP. O fluxo de integração contínua executou \textit{smokes} com PostgreSQL e Kafka reais para os caminhos normais e para uma interrupção do banco. A cobertura desses \textit{smokes} é delimitada na Seção~\ref{sec:validacao-funcional}; a campanha definitiva acrescentou ensaios de carga e falhas, apresentados no Capítulo~\ref{cap:resultados}. Verificações funcionais documentam o software e não respondem isoladamente à pergunta experimental.

O piloto verificou geração de carga, coleta, resolução temporal, limites e duração. Correções instrumentais pertencem a essa fase. O protocolo foi congelado antes da coleta definitiva; a revisão R1 rejeitada permanece preservada. A emenda revisão analítica posterior, detalhada adiante, alterou somente análise offline e controle de admissão, sem mudar instrumentos, imagens, parâmetros ou repetir UUIDs.

\subsection{Cenários de teste}
C1 utilizou carga baixa para observar correção sem saturação. C2 aplicou níveis graduais para avaliar latência e capacidade. C3 produziu uma rajada acima da taxa estável. C4 introduziu indisponibilidade durante entrada ativa. C5 acompanhou a retomada e a drenagem do backlog.

A calibração inicial confirmou uma referência comum de 48 eventos por segundo durante cinco minutos de medição em cada variante, após aquecimento. Essa taxa é uma referência estável observada, não a capacidade máxima. C2 utilizou 12, 24, 36 e 48 eventos por segundo, correspondentes a 25\%, 50\%, 75\% e 100\% da referência; C1 compartilha a condição de 12 eventos por segundo, sem duplicar sua contagem amostral. C3 utilizou a rajada 48--96--48, com duas transições de um segundo. Cada par recebeu a mesma taxa absoluta e conteúdo determinístico correspondente; foram informadas diferenças entre volume programado e efetivamente enviado.

Em C4, a falha escolhida é a interrupção do PostgreSQL nas duas variantes, para comparar a reação à mesma dependência indisponível. A API síncrona pode retornar 503, enquanto a assíncrona ainda pode confirmar publicação com 202 se Kafka estiver disponível; esse aceite não garante persistência posterior. Se o orçamento de tentativas do consumidor se esgotar, a mensagem segue à DLQ e não há reprocessamento automático depois da volta do banco. O piloto fixou duração da interrupção, prazo de drenagem e critério de estabilidade de C5.

Antes de qualquer coleta definitiva, foi autorizada uma revisão do orçamento experimental: o protocolo anterior, de dez repetições e 300 segundos de medição, foi preservado como versão 1, sem reclassificar calibrações ou demonstrações como resultados definitivos. A versão 2 manteve as 13 combinações de condição e variante, com cinco repetições por combinação, totalizando 65 execuções. Cada coleta teve 60 segundos de aquecimento excluídos da análise, intervalo de 12 segundos, 180 segundos de medição e 90 segundos de observação posterior. Os cenários de falha utilizaram 12 eventos por segundo e interrupção solicitada de 40 segundos, iniciada aproximadamente aos 60 segundos da medição; os instantes efetivos foram registrados. Foram disponibilizados 128 usuários virtuais disponíveis nas condições normais e 176 nas falhas, sem pressupor simultaneidade constante.

A redução é uma escolha explícita de orçamento, não um cálculo de poder estatístico nem demonstração de equivalência com o tamanho amostral anterior. Cinco repetições reduzem a precisão entre execuções e a janela menor pode deixar de observar efeitos tardios. A análise apresenta valores por repetição, média, mediana, desvio-padrão, extremos e diferenças pareadas, sem tratar eventos individuais como repetições independentes ou converter ausência de diferença clara em equivalência. Resultados mistos ou insuficientes permanecem inconclusivos.

Na revisão R2, somente a sonda periódica de saúde do broker experimental foi substituída: em vez de iniciar uma JVM administrativa para listar tópicos a cada dez segundos, uma conexão TCP local verifica a disponibilidade de transporte. Essa sonda leve não comprova funcionamento semântico do Kafka; criação efetiva de tópicos e verificação da API continuam obrigatórias antes da carga. O objetivo é não misturar repetidas inicializações da ferramenta administrativa ao processamento medido do broker. Imagens, quotas, regras de negócio, referência de 48 eventos por segundo e critérios de elegibilidade são mantidos. O ensaio R1 que motivou a investigação permanece identificado e não será aprovado retrospectivamente ou contado como definitivo.

O congelamento da versão 2 foi condicionado a nove validações identificadas separadamente: modelo comum nas duas variantes; referência de 48 eventos por segundo e rajada nas duas; interrupção do PostgreSQL nas duas; e interrupção do consumidor apenas na assíncrona. As sete validações de carga utilizaram a janela e os limites propostos para os definitivos. Exigiram integridade instrumental, margem dos coletores, elegibilidade da referência e recuperação sustentada durante entrada ativa nos cenários de falha. Os nove casos foram aprovados e auditados em 27 de outubro de 2026, com preservação dos dados e da tentativa R1 não elegível. Não foram alongadas retrospectivamente as janelas ou afrouxados os critérios; a aprovação libera a campanha, mas não substitui suas 65 repetições.

A ordem das condições foi embaralhada pela semente registrada 20261026, em cinco blocos completos. A primeira variante foi contrabalançada globalmente; em cada condição pareada a proporção foi 3/2 ou 2/3, pois cinco repetições não permitem igualdade exata. Entre execuções, somente a tabela sintética foi reinicializada, com tópico e grupo novos; volumes, broker e caches não foram fisicamente zerados. As evidências anteriores foram exportadas e preservadas antes da próxima execução. Validações, calibrações e demonstrações não integram as 65 repetições definitivas.

\subsection{Métricas coletadas}
A duração HTTP nativa do k6, \texttt{http\_req\_duration}, não inclui todos os custos anteriores de conexão e bloqueio; o intervalo instrumentado entre envio e retorno é uma medida distinta. A comparação principal de resposta utilizou o p95 das respostas válidas esperadas (201 na síncrona e 202 na assíncrona, com UUID correspondente), acompanhado da quantidade de observações e das categorias de todas as respostas, inclusive erros. Na síncrona, o sucesso sucede a persistência; na assíncrona, sucede o ACK de publicação. A latência de conclusão utiliza, em ambas, o início do envio e o marco comum pós-commit, conciliado com o banco. Resposta e conclusão são apresentadas separadamente, evitando comparar aceitação de uma variante com conclusão da outra ou interpretar erros rápidos como melhor desempenho.

Throughput concluído é a quantidade de UUIDs únicos efetivamente persistidos na janela de medição dividida pela duração dessa janela. Eventos concluídos somente durante a drenagem são contabilizados à parte. Entradas oferecidas, respostas aceitas e linhas persistidas constituem conjuntos distintos. Uma confirmação pós-commit cujo registro de envio esteja ausente é uma inconsistência de coleta, e não uma amostra válida de latência.

Taxa de erro foi discriminada em HTTP, publicação, validação, persistência, conflito e falha terminal. O lag de Kafka resulta da diferença entre posição final e posição confirmada por partição. A diferença entre eventos aceitos e persistidos constitui outro indicador e não foi tratada como sinônimo de lag. Tentativas locais de retry, releituras da mesma posição Kafka e reenvios de um mesmo UUID também foram contados separadamente.

Ao fim do prazo de observação, os eventos aceitos sem persistência foram separados em falhas terminais registradas na DLQ, eventos ainda pendentes e ausências não conciliadas. Uma mensagem em DLQ foi retida para inspeção; não representa sucesso de auditoria, mas tampouco demonstra desaparecimento no transporte. A afirmação de perda exige investigar a ausência na persistência e nas evidências de retenção. Submissões de resultado incerto serão conciliadas sem serem automaticamente classificadas como aceites confirmados ou perdas.

Duplicidade persistida corresponde a mais de uma linha para o mesmo UUID e deve ser zero. Conflito corresponde ao mesmo UUID associado a conteúdo diferente. O tempo de recuperação foi medido desde a retomada até o atendimento do critério de estabilidade de backlog e throughput fixado no protocolo. Uma eventual observação isolada de lag zero não substituirá esse critério científico.

\subsection{Emenda analítica após início da campanha}
Em 27 de outubro de 2026, após 37 tentativas definitivas, foi autorizada uma emenda analítica versionada, identificada tecnicamente no Apêndice~\ref{ap:rastreabilidade}. Sua motivação foi a coexistência, em uma coleta síncrona de interrupção do PostgreSQL, de HTTP 503 e registro persistido, mas sem marco pós-commit. A regra original invalidava globalmente a discrepância entre banco e marco. O resumo inválido, o motivo e os dados brutos permanecem preservados; a emenda não declara essa coleta integralmente válida pelo critério original. Sua aplicação é retrospectiva nas primeiras 37 tentativas e declarada antes das 28 restantes, não uma regra predeclarada antes de todos os dados.

A qualificação parcial é restrita a eventos medidos no cenário síncrono de interrupção do PostgreSQL: cliente completo, HTTP 503, conteúdo e hash corretos no banco, um marco de recebimento e um de falha por indisponibilidade do banco, sem confirmação de persistência ou commit, retry, replay, DLQ ou outra inconsistência instrumental. A mesma regra foi aplicada a todo o conjunto definitivo. Dados incompletos de cliente, sucesso sem marco, marco sem registro, conteúdo incorreto, relógio inválido, coleta insuficiente, falta de folga ou casos fora desse escopo continuam bloqueios. Não há nova carga de piloto, mudança nas imagens ou instrumentos, nem repetição ou substituição da observação.

As métricas têm elegibilidade própria. HTTP 503 continua erro, ainda que o registro apareça na exportação final; contagens finais e recursos auditados permanecem utilizáveis. A execução parcial e sua contraparte são excluídas apenas da comparação pareada primária de latência de conclusão. Sua distribuição parcial é apresentada separadamente como sensibilidade, com quantidades de tempos observados e registros persistidos. Nenhum timestamp será imputado: o campo preenchido durante o INSERT, o horário do erro e o comando de retomada não substituem a conclusão da transação.

Para vazão, sejam $c$ os commits observados e conciliados dentro da janela, $u$ os registros medidos persistidos sem marco temporal e $T$ a duração da medição. A coleta parcial tem limites conservadores $[c/T,(c+u)/T]$, não um valor preenchido nem intervalo de confiança. Recuperação será descrita pelos marcos observados, sem reconstruir o primeiro commit real desconhecido. A ausência temporal na fronteira da falha não será assumida aleatória ou inofensiva por ser numericamente pequena.

O tamanho efetivo é informado por indicador. O catálogo final registra 64 coletas completas e uma parcial nas 65 coletas; a comparação pareada de interrupção do PostgreSQL tem quatro pares completos para latência de conclusão, embora indicadores identificados mantenham cinco observações. A análise estrita original exclui o par afetado de todos os indicadores pareados como sensibilidade, sem alterar os resumos ou repor a coleta. As conclusões permanecem limitadas pela decisão pós-hoc, pela cobertura temporal e pelo pequeno número de repetições.

\subsection{Critérios de aceitação}
O núcleo será considerado funcional nos cenários verificados quando: eventos válidos forem persistidos; entradas inválidas forem rejeitadas; 201 ocorrer após confirmação da transação; 202 depender de ACK individual; duplicatas preservarem uma linha; conflitos preservarem o registro original; retry e DLQ obedecerem à política; e offsets de origem não forem confirmados antes da persistência ou retenção na DLQ.

A coleta definitiva será aceita se configuração e snapshot forem registrados, instrumentos estiverem validados, protocolo estiver congelado, aquecimento estiver separado, repetições previstas forem identificadas, ordem estiver registrada e dados brutos permitirem conciliação. Falhas do instrumento, reinicializações externas e dados incompletos permanecem causas de invalidade sob os critérios originais; a revisão analítica acrescenta somente a qualificação parcial e a elegibilidade por métrica explicitadas acima. Um resultado desfavorável à hipótese não invalida a execução.

\subsection{Protocolo de execução e reprodutibilidade}
O executor implementado inicia o projeto Compose isolado com imagens existentes, aguarda prontidão, confirma esquema e tópicos, registra commit e configuração, executa aquecimento e carga, registra falha e retomada, aguarda observação posterior e exporta logs, offsets e registros com checksums. O controlador reduzido acrescenta plano predeclarado, exclusão mútua e diário de progresso, recusando sobrescrever UUIDs existentes ou iniciar uma coleta concorrente. As imagens calibradas não são reconstruídas por esse procedimento.

Cada \texttt{run\_id} possui manifesto, eventos enviados, saída do gerador, logs JSON, estatísticas de recursos, offsets e exportação do banco. O controlador preserva diretórios anteriores e interrompe para investigação se houver falha de coleta ou margem insuficiente do instrumento. O protocolo e os comandos são versionados no repositório. Apenas o projeto interativo deste TCC é parado durante a medição e retomado ao final ou após saída tratável; projetos alheios não são interrompidos. A validade da revisão reduzida e os resultados definitivos foram verificados e são relatados no Capítulo~\ref{cap:resultados}.

\subsection{Ameaças à validade}
A validade interna é afetada por concorrência com outros processos, caches, aquecimento, virtualização, limites de recursos, coleta de telemetria e ordem dos ensaios. A mitigação inclui registrar o ambiente, utilizar configurações estáveis, reduzir processos concorrentes durante a coleta, repetir e alternar a ordem. O custo do broker adicional pertence ao objeto comparado; por isso, a pesquisa não isola universalmente um estilo arquitetural.

A validade de construto depende dos marcos temporais e da classificação dos eventos. Separar ACK, commit e resposta HTTP reduz ambiguidade. Resolução de relógios, atraso da observação e custo de emissão dos logs limitam a precisão. Lag, fila pendente, DLQ e perda são estados distintos, cujo tratamento deve ser verificável nos dados.

A validade de conclusão depende de repetições suficientes e de sua independência. Com cinco repetições, a análise examina valores individuais, dispersão e magnitudes das diferenças pareadas, sem interpretar milhares de eventos correlacionados como milhares de experimentos ou alegar superioridade estatística. A ordem contrabalançada não elimina toda interferência ou dependência do ambiente compartilhado. Resultados de validação funcional não serão extrapolados para desempenho.

A validade externa é limitada por dados sintéticos, domínio simplificado, uma implementação, broker único, versões e ambiente local. Não serão generalizadas conclusões para todas as APIs REST ou arquiteturas EDA. A preservação do protocolo, das versões e dos dados permite replicações que investiguem essas limitações.
